% ==============================================================================
% 04_cadre_legal.tex — UE SIS589
% ==============================================================================
\chapter{Cadre Légal, Normatif et Admissibilité des Preuves Numériques}
\label{chap:cadre-legal}

\epigraph{\itshape « Une preuve mal collectée est une preuve perdue. »}
         {--- Principe ACPO, \textit{Good Practice Guide for Digital Evidence}}

\begin{boxobjectifs}
\begin{enumerate}
  \item Connaître les normes internationales structurant la réponse aux incidents
        et l'investigation numérique (ISO~27035, ISO~27037, RFC~3227, NIST~SP~800-86).
  \item Appliquer le cadre juridique camerounais (Loi~2010/012, Loi~2010/013)
        aux situations professionnelles courantes.
  \item Comprendre les conditions d'admissibilité des preuves numériques
        devant les juridictions.
  \item Identifier les obligations de notification et de coopération avec les
        autorités (ANTIC, parquet).
  \item Articuler le droit africain (Convention de Malabo) avec les pratiques
        d'investigation.
\end{enumerate}
\end{boxobjectifs}

% ==============================================================================
\section{Normes internationales}
\label{sec:normes-int}
% ==============================================================================

\subsection{ISO/IEC~27035~: Gestion des incidents de sécurité}

\begin{boxdef}
\textbf{ISO/IEC~27035~:} Norme internationale décrivant les principes et
processus de gestion des incidents de sécurité de l'information. Structurée
en trois parties~: (1)~principes, (2)~lignes directrices de planification et
préparation, (3)~lignes directrices de réponse aux incidents.
\end{boxdef}

Le cycle ISO~27035 comprend cinq phases~:
\begin{enumerate}
  \item \textbf{Planification et préparation~:} politique de gestion des incidents,
        constitution du CSIRT, formation.
  \item \textbf{Détection et signalement~:} identification des événements,
        qualification, escalade.
  \item \textbf{Évaluation et décision~:} confirmation de l'incident, évaluation
        de l'impact, décision de réponse.
  \item \textbf{Réponse~:} confinement, éradication, rétablissement, investigation
        forensique.
  \item \textbf{Leçons tirées~:} RETEX, amélioration du dispositif.
\end{enumerate}

\subsection{ISO/IEC~27037~: Collecte de preuves numériques}

\begin{boxdef}
\textbf{ISO/IEC~27037~:} Lignes directrices pour l'identification, la collecte,
l'acquisition et la préservation des preuves numériques. Applicable par les
enquêteurs, les CSIRT et les équipes forensiques. Principe fondateur~:
\textbf{l'intégrité de la preuve est prioritaire sur la rapidité de la collecte}.
\end{boxdef}

Quatre exigences fondamentales~:
\begin{itemize}
  \item \textbf{Admissibilité~:} la preuve peut être produite et acceptée
        devant un tribunal.
  \item \textbf{Authenticité~:} la preuve est ce qu'elle prétend être
        (pas altérée, pas fabriquée).
  \item \textbf{Complétude~:} toutes les preuves pertinentes ont été collectées.
  \item \textbf{Fiabilité~:} les méthodes utilisées sont reproductibles et
        documentées.
\end{itemize}

\subsection{RFC~3227~: Ordre de collecte des preuves}

\begin{boxdef}
\textbf{RFC~3227 (\textit{Guidelines for Evidence Collection and Archiving})~:}
Document de meilleures pratiques (BCP~55) définissant l'ordre de collecte
des preuves numériques en fonction de leur volatilité. Principe~: collecter
d'abord ce qui disparaît le plus vite.
\end{boxdef}

\begin{table}[H]
\centering\small
\caption{Ordre de volatilité selon RFC~3227}
\label{tab:ordre-volatilite}
\rowcolors{2}{TableauPair}{TableauImpair}
\begin{tabular}{p{0.8cm}p{5cm}p{7cm}}
\toprule
\tableauHeader{\#} & \tableauHeader{Source} &
\tableauHeader{Durée de vie / Remarque} \\
\midrule
1 & Registres CPU, cache L1/L2 & Quelques nanosecondes --- impossible à capturer \\
2 & Mémoire vive (RAM) & Perdue à l'extinction. Dumper avant tout arrêt \\
3 & État réseau (tables ARP, connexions) & Quelques secondes à minutes \\
4 & Processus en cours & Durée de vie variable \\
5 & Disque dur / SSD & Persistant mais modifiable (accès OS) \\
6 & Journaux système distants & Semi-persistant (rotation) \\
7 & Données hors-ligne (backups, archives) & Persistant \\
8 & Configuration physique / topologie & Quasi-permanent \\
\bottomrule
\end{tabular}
\end{table}

\subsection{NIST SP~800-86~: Intégration de la forensique dans l'IR}

\begin{boxdef}
\textbf{NIST SP~800-86~:} Guide d'intégration des techniques forensiques
dans le processus de réponse aux incidents. Définit quatre phases
de l'investigation~: collecte, examen, analyse, rapport.
\end{boxdef}

\begin{table}[H]
\centering\small
\caption{Phases d'investigation selon NIST SP~800-86}
\label{tab:nist-phases}
\rowcolors{2}{TableauPair}{TableauImpair}
\begin{tabular}{p{2.5cm}p{5cm}p{6cm}}
\toprule
\tableauHeader{Phase} & \tableauHeader{Activités} &
\tableauHeader{Livrables} \\
\midrule
Collecte &
Identification et acquisition des preuves~: images disques, dumps RAM,
captures réseau, journaux &
PV de collecte, hashes d'intégrité \\
Examen &
Extraction des données pertinentes depuis les images~: fichiers, métadonnées,
données supprimées &
Fichiers extraits, liste d'artefacts \\
Analyse &
Reconstruction de la timeline, corrélation, identification des TTPs,
attribution &
Timeline annotée, rapport d'analyse \\
Rapport &
Documentation des découvertes, conclusions, recommandations,
forme adaptée au destinataire &
Rapport technique, résumé direction \\
\bottomrule
\end{tabular}
\end{table}

\subsection{Principes ACPO}

Les quatre principes ACPO (\textit{Association of Chief Police Officers},
Royaume-Uni) sont devenus un standard de facto pour la collecte de
preuves numériques~:

\begin{enumerate}
  \item \textbf{Principe~1~:} Aucune action ne doit modifier les données
        sur un support qui pourrait être ultérieurement présenté en justice.
  \item \textbf{Principe~2~:} Si une personne estime nécessaire d'accéder
        aux données originales, cette personne doit être compétente pour
        le faire et capable d'expliquer la pertinence et les implications
        de ses actions.
  \item \textbf{Principe~3~:} Une piste d'audit complète de toutes les
        actions effectuées sur les preuves numériques doit être créée et
        préservée.
  \item \textbf{Principe~4~:} La personne responsable de l'investigation
        est responsable du respect de ces principes.
\end{enumerate}

% ==============================================================================
\section{Cadre juridique camerounais}
\label{sec:droit-cameroun}
% ==============================================================================

\subsection{Loi n°~2010/012 relative à la cybersécurité et à la cybercriminalité}

\begin{boxdef}
\textbf{Loi~2010/012 du 21~décembre~2010~:} Texte fondateur du droit
de la cybersécurité au Cameroun. Elle définit les infractions liées
aux systèmes d'information, les obligations des opérateurs, les pouvoirs
des enquêteurs et le cadre de la coopération internationale.
\end{boxdef}

\textbf{Infractions principales définies~:}
\begin{itemize}
  \item Accès frauduleux à un système d'information (Art.~62)~:
        peine jusqu'à 2~ans d'emprisonnement et 5~millions~FCFA d'amende.
  \item Maintien frauduleux dans un système (Art.~63).
  \item Atteinte à l'intégrité d'un système ou de données (Art.~64--66).
  \item Interception illégale de communications (Art.~67).
  \item Falsification de données (Art.~68).
\end{itemize}

\textbf{Obligations des opérateurs (Art.~44--50)~:}
\begin{itemize}
  \item Conservation des données de connexion pendant une durée
        déterminée par décret.
  \item Coopération avec les autorités judiciaires sur réquisition.
  \item Signalement des violations de données à l'ANTIC.
\end{itemize}

\textbf{Pouvoirs des enquêteurs~:}
\begin{itemize}
  \item Saisie de matériels informatiques sur autorisation judiciaire (Art.~83).
  \item Accès aux données chiffrées~: obligation de remise des clés sur
        réquisition (Art.~85).
  \item Interception légale de communications électroniques (Art.~87--89),
        sous contrôle du parquet.
\end{itemize}

\subsection{Loi n°~2010/013 régissant les communications électroniques}

Cette loi complète la~2010/012 en définissant le cadre des opérateurs de
télécommunications, les obligations de couverture, de confidentialité
des communications et les droits des usagers.

Elle est pertinente pour la cybersécurité dans son volet~: obligation de
sécurisation des infrastructures (Art.~79), responsabilité des prestataires
hébergeurs, et encadrement de l'interception légale.

\subsection{L'ANTIC~: rôle et pouvoirs}

\begin{boxdef}
\textbf{ANTIC (\textit{Agence Nationale des Technologies de l'Information
et de la Communication})~:} Autorité de régulation et de supervision de
la cybersécurité au Cameroun. Ses missions incluent la certification des
équipements cryptographiques, la surveillance des réseaux publics,
le traitement des incidents de cybersécurité signalés et la coopération
internationale.
\end{boxdef}

En cas d'incident, l'ANTIC peut~:
\begin{itemize}
  \item Être notifiée et intervenir techniquement.
  \item Requérir la communication de journaux et données techniques.
  \item Orienter les entreprises vers les autorités judiciaires compétentes.
\end{itemize}

% ==============================================================================
\section{Convention de Malabo et cadre africain}
\label{sec:malabo}
% ==============================================================================

\begin{boxdef}
\textbf{Convention de l'Union Africaine sur la Cybersécurité et la Protection
des Données Personnelles (Convention de Malabo, 2014)~:} Traité continental
définissant un cadre harmonisé pour la cybersécurité, la cybercriminalité,
la protection des données et les transactions électroniques en Afrique.
Le Cameroun est signataire.
\end{boxdef}

Quatre axes principaux~:
\begin{enumerate}
  \item \textbf{Transactions électroniques~:} valeur juridique des documents
        et signatures électroniques.
  \item \textbf{Protection des données personnelles~:} droits des personnes,
        obligations des responsables de traitement (parallèle au RGPD).
  \item \textbf{Sécurité des SI~:} obligations de sécurisation, mécanismes
        de certification, CSIRT nationaux.
  \item \textbf{Cybercriminalité~:} incrimination des infractions
        informatiques, coopération judiciaire transfrontalière.
\end{enumerate}

\begin{boiteRemarque}{Impact pratique pour les entreprises}
La Convention de Malabo impose aux États membres de disposer d'une autorité
nationale de protection des données et d'un CSIRT national. Pour les entreprises
opérant dans plusieurs pays africains, elle constitue le socle de la
coopération en cas d'incident transfrontalier.
\end{boiteRemarque}

% ==============================================================================
\section{Admissibilité des preuves numériques}
\label{sec:admissibilite}
% ==============================================================================

\subsection{Conditions générales d'admissibilité}

Pour qu'une preuve numérique soit recevable devant une juridiction, elle
doit satisfaire aux conditions suivantes~:

\begin{description}
  \item[Légalité de la collecte.] La preuve doit avoir été obtenue dans
        le respect des droits (sans violation de vie privée non autorisée,
        sans accès non autorisé). Une preuve obtenue illégalement est
        irrecevable.
  \item[Authenticité.] La preuve est ce qu'elle prétend être. Les hashes
        cryptographiques (SHA-256) calculés au moment de la saisie,
        conservés dans le PV de collecte, garantissent qu'elle n'a pas
        été altérée.
  \item[Intégrité.] La preuve n'a pas été modifiée depuis sa collecte.
        La chaîne de custody documente chaque accès.
  \item[Fiabilité de la méthode.} Les outils et méthodes utilisés sont
        reconnus et documentés (FTK Imager, dd, Autopsy, Volatility).
\end{description}

\subsection{La chaîne de custody}

\begin{boxdef}
\textbf{Chaîne de custody (\textit{Chain of Custody})~:} Document traçant
l'historique complet d'une pièce à conviction numérique depuis sa collecte
initiale jusqu'à sa présentation en justice~: qui l'a manipulée, quand,
pourquoi, et dans quel état elle se trouvait à chaque étape.
Toute rupture de la chaîne fragilise l'admissibilité de la preuve.
\end{boxdef}

\begin{boxcustody}
\textbf{Modèle de chaîne de custody --- LiZbiZ-SIS}

\begin{itemize}
  \item \textbf{Pièce~:} Image disque \texttt{srv-web-01-sda.dd}
  \item \textbf{Date/heure de collecte~:} 12/04/2026, 04:22~UTC
  \item \textbf{Collecteur~:} MINKA T.E., Analyste N2 SOC LiZbiZ
  \item \textbf{Méthode~:} FTK Imager v4.7.1.2, copie bit à bit
  \item \textbf{SHA-256~:} \hash{sha256}{3a7d4f8b...c2e91}
  \item \textbf{Support de stockage~:} Disque externe chiffré
        (\texttt{VeraCrypt}, SN: XYZ-20260412)
  \item \textbf{Conservation~:} Coffre-fort salle serveur LiZbiZ,
        scellé n°~2026-04-001
  \item \textbf{Accès ultérieur~:} 14/04/2026, 09:00~UTC ---
        Expertise judiciaire, tribunal de Yaoundé, juge FOFANA
\end{itemize}
\end{boxcustody}

\subsection{Valeur probatoire des logs SIEM}

Un journal exporté en temps réel vers un SIEM tiers dispose d'une valeur
probatoire plus forte qu'un journal stocké localement, car~:
\begin{itemize}
  \item L'horodatage est fourni par le serveur SIEM (externe à la machine
        compromise), difficile à falsifier rétrospectivement.
  \item L'intégrité de la transmission peut être attestée par les
        certificats TLS et les logs du collecteur.
  \item Le SIEM constitue un tiers de confiance opérationnel.
\end{itemize}

% ==============================================================================
\section{Obligations de notification}
\label{sec:notification}
% ==============================================================================

\begin{table}[H]
\centering\small
\caption{Délais et destinataires de notification selon les référentiels}
\label{tab:notification}
\rowcolors{2}{TableauPair}{TableauImpair}
\begin{tabular}{p{4cm}p{2.5cm}p{3cm}p{4cm}}
\toprule
\tableauHeader{Référentiel} & \tableauHeader{Délai} &
\tableauHeader{Destinataire} & \tableauHeader{Périmètre} \\
\midrule
RGPD Art.~33 (UE) & 72~heures & Autorité de contrôle (CNIL) &
Violation données personnelles \\
Directive NIS2 (UE) & 24h (alerte) + 72h (rapport) & CSIRT national + ENISA &
Opérateurs services essentiels \\
Loi~2010/012 Cameroun & Non précisé & ANTIC, parquet & Incidents majeurs \\
PCI-DSS (Art.~12.10) & Immédiat & Banques acquéreuses & Compromission données cartes \\
ISO~27035 & Défini par politique interne & CSIRT interne, management & Tous incidents confirmés \\
\bottomrule
\end{tabular}
\end{table}

\begin{boxlizbiz}
\textbf{LiZbiZ-SIS~: obligations applicables}

LiZbiZ-SIS traite des données de paiement mobile money (MTN MoMo,
Orange Money)~: elle est soumise aux exigences PCI-DSS si elle stocke,
traite ou transmet des données de cartes. Elle est également soumise à
la Loi~2010/012 pour la conservation des journaux de connexion et la
coopération avec l'ANTIC. En cas de fuite de données clients, elle doit
notifier l'ANTIC sans délai injustifié, documenter la violation et
prendre les mesures de remédiation immédiates.
\end{boxlizbiz}

% ==============================================================================
\section*{Résumé}
% ==============================================================================

\begin{boxresume}
\begin{itemize}
  \item ISO~27035~: cycle de gestion des incidents en 5~phases.
        ISO~27037~: collecte de preuves (admissibilité, authenticité,
        complétude, fiabilité).
  \item RFC~3227~: ordre de volatilité (RAM~$\to$~disque~$\to$~archives).
        NIST~SP~800-86~: collecte, examen, analyse, rapport.
  \item Principes ACPO~: ne pas altérer les données originales~;
        documenter tout accès~; maintenir la chaîne de custody.
  \item Loi~2010/012~: infractions informatiques, pouvoirs des
        enquêteurs, obligations des opérateurs. ANTIC~: autorité compétente.
  \item Convention de Malabo~: cadre africain harmonisé (données
        personnelles, cybercriminalité, coopération transfrontalière).
  \item Chaîne de custody~: traçabilité complète de la preuve.
        Logs SIEM distants~: valeur probatoire supérieure.
\end{itemize}
\end{boxresume}

\begin{boxexo}
\textbf{Ex.~4.1 — Qualification juridique.}
Un technicien LiZbiZ découvre que son collègue a accédé à la base de
données clients en dehors de ses attributions pour exporter 50~000 dossiers.
Qualifiez les infractions potentielles selon la Loi~2010/012. Quelles
preuves faut-il collecter~? Quel est le délai pour notifier l'ANTIC~?

\medskip\textbf{Ex.~4.2 — Chaîne de custody.}
Vous devez saisir un laptop Windows suspect sur le bureau d'un employé.
Rédigez le PV de saisie complet (contexte, méthode, hashes, stockage)
en appliquant les principes ACPO et ISO~27037.

\medskip\textbf{Ex.~4.3 — Politique de rétention.}
LiZbiZ-SIS est en négociation avec un client nigérian et un client
ivoirien. Identifiez les référentiels applicables à chacun et proposez
une politique de rétention unifiée compatible avec les trois juridictions
(Cameroun, Nigeria, Côte d'Ivoire).
\end{boxexo}
